\clearpage
\chapter*{Erratum à l'exposé n\textsuperscript{o} 149}
\addcontentsline{toc}{chapter}{Erratum à l'exposé n\textsuperscript{o} 149}
\setcounter{footnote}{0}
\src{298}
\begingroup
\makeatletter
\renewcommand\@makefnmark{\hbox{\normalfont(\textsuperscript{\@thefnmark})}}
\makeatother

\noindent 1. \uwave{Théorème de dualité pour les faisceaux algébriques
cohérents}.

\begin{center}
[Séminaire Bourbaki, t. 9, 1956/57, n\textsuperscript{o} 149, 25 p.]
\end{center}

\noindent\uline{Remarque sur la page 14}. - Comme je l'ai signalé dans ma
conférence au Congrès international des Mathématiciens en 1958 \footnote{
GROTHENDIECK (Alexander). - The cohomology theory of abstract algebraic
varieties, Proceedings of the international Congress of Mathematicians
[1958. Edinburgh] ; p. 103-118. - Cambridge, at the University Press, 1960.},
les questions soulevées ici sont maintenant complètement résolues.

Le lecteur trouvera d'autres renseignements sur la dualité des faisceaux
cohérents dans loco citato \footnotemark[1], p. 112-115, dans \uwave{EGA} III,
2e partie (en préparation), et, en attendant, dans \uwave{SGA} 1962.
Un traitement plus systématique se trouvera dans un chapitre ultérieur de \uwave{EGA}
(chapitre IX dans le plan prévu).
\endgroup
